\section{Basic Concepts of Acoustics}
Since we want to take a closer look at the origin of noises, sounds, and music at the Junior Academy, we will first familiarize ourselves with the basic concepts of acoustics. The field of acoustics primarily deals with the \textbf{study of sound} and its propagation. There are three main branches of acoustics:
\begin{itemize}
    \item \textbf{physical acoustics}: physical acoustics deals with explaining acoustic phenomena using classical mechanics
    \item \textbf{physiological acoustics}: physiological acoustics studies the reception and transmission of sound within the human auditory system
    \item \textbf{psychological acoustics}: psychological acoustics investigates how acoustic nerve stimuli are translated into auditory perception, i.e., it explains why sounds and tones are perceived the way they are
\end{itemize}
In the following, we will mainly focus on physical acoustics.

\subsection{Sound}
In science, the term ``sound'' describes mechanical vibrations in an elastic medium (e.g., air or water). You may also have heard the term ``sound wave.'' You can imagine a \textbf{sound wave} similarly to a wave in water. When air particles move, i.e., oscillate back and forth, they transfer their motion to neighboring air particles. This \textbf{motion of air particles}, however, is not visible in contrast to water waves. Later, we will take a closer look at sound waves.
 
\subsection{Tone}
You are probably familiar with the term tone as a ``sound event'' produced by a musical instrument or the human voice. From a physical perspective, a tone is determined by the frequency of a harmonic sound wave. The frequency indicates how many oscillations the air particles perform per second and is therefore crucial for which tone you perceive. A tone with a higher frequency is perceived as a ``higher'' tone. Thus, the physical concept of tone describes what you likely know as \textbf{pitch}.

\subsection{Timbre}
The physical concept of sound (timbre) corresponds roughly to what is commonly referred to as ``tone'' in music and everyday language. For example, if you play a tone with the same frequency on a piano and a violin, you can still distinguish which instrument produced the sound. This is because what we hear is not just a single tone but a \textbf{superposition of sound waves with different frequencies}. This combination of overlapping sound waves is called a sound (timbre). The reason you can recognize an instrument by hearing a melody is that each instrument has a characteristic composition of different frequencies. Without these differences in sound, music would be very monotonous.

%%% Add some more motivation for the following chapters.